% Physical pages: 412--417 (printed pages 389--394).
% Figures: none.
% Uncertainties: none.

\subsection{Lagrangian Version of the Path-Integral Formula}

The integrand in the exponential in Eqs. \eqref{eq:9.1.38} or \eqref{eq:9.2.17} looks like
the Lagrangian $L$ associated with the Hamiltonian $H$. This appearance is
somewhat misleading because here the `momenta' $p_a(t)$ or $p_n(\mathbf{x},t)$
are independent variables, not yet related to $q_a(t)$ or $q_n(\mathbf{x},t)$
or their derivatives. However, there is a large and important class of
theories in which the integral over the `momenta' can be done by just replacing
them with the values dictated by the canonical formalism, in which case the
integrand in the exponential in the path integrals really is the Lagrangian.

These are theories with a Hamiltonian that is quadratic in the `momenta'---in
the language of field theory
\begin{align}
H[Q,P]={}&\frac12\sum_{nm}\int d^3x\,d^3y\,
 A_{xn,ym}[Q]P_n(\mathbf{x})P_m(\mathbf{y})
\notag\\
&+\sum_n\int d^3x\,B_{xn}[Q]P_n(\mathbf{x})+C[Q]
\tag{9.3.1}\label{eq:9.3.1}
\end{align}
with a `matrix' $A$ that is real, symmetric, positive, and non-singular. The
argument of the exponential in Eq. \eqref{eq:9.2.17} is then quadratic in the $p$s:
\begin{align}
&\int d\tau\left\{\int d^3x\sum_n p_n(\mathbf{x},\tau)
 \dot q_n(\mathbf{x},\tau)-H[q(\tau),p(\tau)]\right\}
\notag\\
&\quad=-\frac12\sum_{nm}\int d^3x\,d^3y\,d\tau\,d\tau'\,
 \mathcal A_{xn\tau,ym\tau'}[q]p_n(\mathbf{x},\tau)p_m(\mathbf{y},\tau')
\notag\\
&\qquad-\sum_n\int d^3x\int d\tau\,
 \mathcal B_{xn\tau}[q]p_n(\mathbf{x},\tau)-\mathcal C[q],
\tag{9.3.2}\label{eq:9.3.2}
\end{align}
where
\begin{align}
\mathcal A_{xn\tau,ym\tau'}[q]
&\equiv A_{xn,ym}[q(\tau)]\delta(\tau-\tau'),
\tag{9.3.3}\label{eq:9.3.3}\\
\mathcal B_{xn\tau}[q]
&\equiv B_{xn}[q(\tau)]-\dot q_n(\mathbf{x},\tau),
\tag{9.3.4}\label{eq:9.3.4}\\
\mathcal C[q]&\equiv\int d\tau\,C[q(\tau)].
\tag{9.3.5}\label{eq:9.3.5}
\end{align}
Now, in general the integral of the exponential of a quadratic expression
like \eqref{eq:9.3.2} will be proportional to the exponential evaluated at the
stationary point of its argument. For a finite number of real variables
$\xi_s$, this formula reads
\begin{align}
&\int_{-\infty}^{\infty}\left(\prod_s d\xi_s\right)
 \exp\left\{-\frac{i}{2}\sum_{sr}\mathcal A_{sr}\xi_s\xi_r
 -i\sum_s\mathcal B_s\xi_s-i\mathcal C\right\}
\notag\\
&\qquad=
 \left(\operatorname{Det}\left[\frac{i\mathcal A}{2\pi}\right]\right)^{-1/2}
 \exp\left\{-\frac{i}{2}\sum_{sr}\mathcal A_{sr}\bar\xi_s\bar\xi_r
 -i\sum_s\mathcal B_s\bar\xi_s-i\mathcal C\right\},
\tag{9.3.6}\label{eq:9.3.6}
\end{align}
where $\bar\xi$ is the stationary point
\begin{equation}
\bar\xi_s=-\sum_r(\mathcal A^{-1})_{sr}\mathcal B_r.
\tag{9.3.7}\label{eq:9.3.7}
\end{equation}
(For a proof of this formula, see the Appendix to this chapter.) Hence, as
long as the $\mathcal O_A,\mathcal O_B$, etc. in Eq. \eqref{eq:9.2.17} are independent
of the $p$s, for such a Hamiltonian we can evaluate the path integral over the
$p$s in Eq. \eqref{eq:9.2.17} by setting these variables at the stationary point of the
quadratic expression in the argument of the exponential. But the variational
derivative of this quadratic is
\begin{align*}
&\frac{\delta}{\delta p_n(\mathbf{x},\tau)}
 \int_{-\infty}^{\infty}d\tau
 \left\{\int d^3x\sum_m\dot q_m(\mathbf{x},\tau)p_m(\mathbf{x},\tau)
 -H[q(\tau),p(\tau)]+i\epsilon\text{ terms}\right\}
\\
&\qquad=\dot q_n(\mathbf{x},\tau)
 -\frac{\delta}{\delta p_n(\mathbf{x},\tau)}H[q(\tau),p(\tau)].
\end{align*}
(The $i\epsilon$ terms depend only on the $q$s.) Thus the stationary `point'
$\bar p_n(\mathbf{x},t)$ where this vanishes is just the value of
$p_n(\mathbf{x},t)$ dictated by the canonical formula
\begin{equation}
\dot q_n(\mathbf{x},t)
=\left[\frac{\delta H[q(t),p(t)]}{\delta p_n(\mathbf{x},t)}\right]_{p=\bar p}.
\tag{9.3.8}\label{eq:9.3.8}
\end{equation}
With $p_n(\mathbf{x},t)$ set equal to this value, the argument of the
exponential in Eq. \eqref{eq:9.2.17} is the ordinary Lagrangian
\begin{equation}
L[q(t),\dot q(t)]\equiv\int d^3x\left(
 \sum_n\dot q_n(\mathbf{x},t)\bar p_n(\mathbf{x},t)
 -H[q(t),\bar p(t)]\right)
\tag{9.3.9}\label{eq:9.3.9}
\end{equation}
and we can write Eq. \eqref{eq:9.2.17} as
\begin{align}
&\bra{\Omega;\mathrm{out}}T\left\{
 \mathcal O_A[Q(t_A)],\mathcal O_B[Q(t_B)],\ldots\right\}
 \ket{\Omega;\mathrm{in}}
\notag\\
&\quad=|\mathcal N|^2\int\prod_{\tau,\mathbf{x},n}dq_n(\mathbf{x},\tau)
 \left(\operatorname{Det}[2i\pi\mathcal A[q]]\right)^{-1/2}
\notag\\
&\qquad\cdot\mathcal O_A[q(t_A)]\mathcal O_B[q(t_B)]\cdots
\notag\\
&\qquad\cdot\exp\left[i\int_{-\infty}^{\infty}d\tau
 \left\{L[q(\tau),\dot q(\tau)]+i\epsilon\text{ terms}\right\}\right].
\tag{9.3.10}\label{eq:9.3.10}
\end{align}
(We have combined the $1/2\pi$ factors in the integrals over the $p_n$ with
the determinant coming from Eq. \eqref{eq:9.3.6}.) This is the desired Lagrangian form
of the path-integral formula.

In deriving Eq. \eqref{eq:9.3.10}, it was necessary to assume that the operators
$\mathcal O_A,\mathcal O_B,\ldots$ were independent of the canonical
`momenta'. This is not as restrictive as it may seem. For instance, in a
scalar field theory for which the canonical conjugate to $\Phi$ is
$\Pi=\dot\Phi$, it is possible to calculate the matrix element of a
time-ordered product of operators, one of which is $\dot\Phi(t)$, by taking
the difference of matrix elements in which this operator is replaced with
$\Phi(t+d\tau)$ and $\Phi(t)$, and then dividing by $d\tau$, with
$d\tau\to0$. Equivalently, as long as $t$ is not equal to any of the other
time-arguments of the operator in Eq. \eqref{eq:9.3.10}, we can simply differentiate
Eq. \eqref{eq:9.3.10} with respect to $t$.

The one serious remaining complication in Eq. \eqref{eq:9.3.10} is the determinant of
$\mathcal A[q]$. If $\mathcal A[q]$ is field-independent, then this is no
problem; we have already noted that overall constants make no contribution to
the connected parts of vacuum expectation values, in which we divide by a
vacuum-vacuum amplitude proportional to the same constant factor. This is the
case for instance for the theory of a set of scalar fields $\Phi_n$ with
non-derivative coupling to each other and/or derivative coupling to external
currents $J_n$. The Lagrangian density here is
\[
\mathcal L=-\sum_n\left[\frac12\partial_\mu\Phi_n\partial^\mu\Phi_n
 +J_n^\mu\partial_\mu\Phi_n\right]-V(\Phi).
\]
An obvious extension of the results of Section 7.5 from one to several
derivatively coupled scalars shows that this Lagrangian implies the Hamiltonian
\[
H=\int d^3x\sum_n\left[\frac12\Pi_n^2+\frac12(\boldsymbol\nabla\Phi_n)^2
 +\mathbf J_n\cdot\boldsymbol\nabla\Phi_n+J_n^0\Pi_n
 +\frac12(J_n^0)^2\right]+\int d^3x\,V(\Phi).
\]
(The $\Phi_n$ are taken to be real scalars, but complex scalars can be
accommodated by separating them into real and imaginary parts.) In general
there is a non-trivial term that is linear in the $\Pi_n$, but the coefficient
of the quadratic term is a constant, just the unit `matrix':
\[
\mathcal A_{xn,x'n'}=\delta^4(x-x')\delta_{nn'}.
\]
The factor $\bigl(\operatorname{Det}[4i\pi\mathcal A[q]]\bigr)^{-1/2}$ in
Eq. \eqref{eq:9.3.10} is here a field-independent constant, and therefore is without
effect.

However, matters are not always so simple. As a second example, let us
consider the so-called non-linear $\sigma$-model, with Lagrangian density
\[
\mathcal L=-\frac12\sum_{nm}\partial_\mu\Phi_n\partial^\mu\Phi_m
 [\delta_{nm}+U_{nm}(\Phi)]-V(\Phi).
\]
A straightforward calculation gives the Hamiltonian as
\[
H=\int d^3x\left[\frac12\Pi_n(\mathbf{1}+U(\Phi))^{-1}_{nm}\Pi_m
 +\frac12\boldsymbol\nabla\Phi_n\cdot\boldsymbol\nabla\Phi_m
 (\mathbf{1}+U(\Phi))_{nm}+V(\Phi)\right].
\]
Here $\mathcal A$ is the field-dependent quantity
\[
\mathcal A_{nx,my}=[\mathbf{1}+U(\Phi(x))]^{-1}_{nm}\delta^4(x-y).
\]
In cases of this sort, the determinant may be reexpressed as a contribution to
the effective Lagrangian, using the relation
$\operatorname{Det}\mathcal A=\exp\operatorname{Tr}\ln\mathcal A$. By
replacing the continuum of spacetime positions with a discrete lattice of
points surrounded by separate regions of very small spacetime volume $\Omega$,
we may interpret the delta function in $\mathcal A_{nx,my}$ as
$\delta^4(x-y)=\Omega^{-1}\delta_{x,y}$, so that
\[
(\ln\mathcal A)_{nx,my}=\delta_{x,y}
 \left[\ln(\mathbf{1}+U(\Phi(x)))-\mathbf{1}\cdot\ln\Omega\right]_{nm}
\]
with the logarithm of a matrix defined now by its power series expansion
\[
\ln(\mathbf{1}+U)=U-\frac{U^2}{2}+\frac{U^3}{3}-\cdots.
\]
To evaluate the trace, we note that
$\sum_x\cdots=\Omega^{-1}\int d^4x\cdots$. The determinant factor here is
then
\[
\operatorname{Det}\mathcal A\propto
 \exp\left[-\Omega^{-1}\int d^4x\,
 \operatorname{tr}\ln[\mathbf{1}+U(\Phi(x))]\right],
\]
where `tr' is to be understood as the trace in an ordinary matrix sense. The
constant of proportionality (which arises from the $-\ln\Omega$ term) is
field-independent and therefore of no present interest. We can regard this
determinant as providing a correction to the effective Lagrangian density
\[
\Delta\mathcal L=-\frac{i}{2}\Omega^{-1}
 \operatorname{tr}\ln[\mathbf{1}+U(\Phi(x))].
\]
The factor $\Omega^{-1}$ may be written as an ultraviolet divergent integral
\[
\Omega^{-1}=\delta^4(x-x)=(2\pi)^{-4}\int d^4p\cdot1.
\]
We shall not show this here, but the extra terms in the Feynman diagrams for
this theory contributed by $\Delta\mathcal L$ could also have been derived in
the canonical formalism by taking account of the equal-time-commutator terms
in the propagator of time-derivatives of the scalar field~\hyperref[ref:9.7]{[7]}. Ignoring this
correction would lead to a spurious dependence of the $S$-matrix on the way
that the scalar field is defined, and would also be inconsistent with any
symmetries of the Lagrangian under transformations of the scalar fields.

\begin{sloppypar}
Even where the factor $(\operatorname{Det}\mathcal A)^{-1/2}$ in the
path-integral formula \eqref{eq:9.3.10} is field-independent, the Lagrangian in this
formula may not be the same as the one with which we started. As an example,
let's consider the theory of a set of real vector fields, with Lagrangian
density
\end{sloppypar}
\begin{align*}
\mathcal L=-\sum_n\biggl[&\frac14(\partial_\mu A_{n\nu}-\partial_\nu A_{n\mu})
 (\partial^\mu A_n^\nu-\partial^\nu A_n^\mu)
\\
&+\frac12m_n^2A_{n\mu}A_n^\mu+J_n^\mu A_{n\mu}\biggr],
\end{align*}
where the currents $J_n^\mu$ are either externally produced c-number
quantities or depend on other fields (in which case terms describing these
other fields are to be added to the Lagrangian). By a simple extension of the
results of Section 7.5, we see that the Hamiltonian is
\begin{align*}
H=\int d^3x\sum_n\biggl[&\frac12\boldsymbol\Pi_n^2
 +\frac12(\boldsymbol\nabla\cdot\mathbf A_n)^2
 +\frac12m_n^2\mathbf A_n^2
 +\frac{1}{2m_n^2}(\boldsymbol\nabla\cdot\boldsymbol\Pi_n)^2
\\
&+\mathbf J_n\cdot\mathbf A_n
 -\frac{1}{m_n^2}J_n^0\boldsymbol\nabla\cdot\boldsymbol\Pi_n
 +\frac{1}{2m_n^2}(J_n^0)^2\biggr],
\end{align*}
again with the understanding that other terms must be added involving any
fields that appear in $J_n^\mu$. Here the coefficient of the quadratic term is
somewhat more complicated than in our first example:
\[
\mathcal A_{nix,mjy}=\delta_{nm}\left[\delta_{ij}\delta^4(x-y)
 -\frac{1}{2m_n^2}\nabla_i\nabla_j\delta^4(x-y)\right],
\]
but it is field-independent, so that the factor
$(\operatorname{Det}\mathcal A)^{-1/2}$ has no effect. On the other hand, the
Lagrangian \eqref{eq:9.3.9} is here not the one with which we started; it is expressed
entirely in terms of $\mathbf A$ and its spacetime derivatives, with no
dependence on any time-component $A^0$. For this reason, the Lorentz invariance
of Eq. \eqref{eq:9.3.10} is far from obvious.

To remedy this, we may reintroduce the auxiliary field. Suppose we add to the
Hamiltonian a term
\[
\Delta H=-\frac12\sum_n m_n^2\int d^3x
 \left[A_n^0-m_n^{-2}\boldsymbol\nabla\cdot\boldsymbol\Pi_n
 +m_n^{-2}J_n^0\right]^2
\]
and integrate over the $A_n^0$ as well as the $\mathbf A_n$ and
$\boldsymbol\Pi_n$. This can only introduce a field-independent overall
factor, since $\Delta H$ is a quadratic in $A^0$ (with a field-independent
coefficient in the term of second order in $A^0$) whose stationary value
vanishes. However, suppose that we now integrate over the $\boldsymbol\Pi_n$
\emph{before} integrating over the $A_n^0$. The Hamiltonian in the path
integral \eqref{eq:9.2.17} is here replaced with
\begin{align*}
H+\Delta H=\int d^3x\sum_n\biggl[&\frac12\boldsymbol\Pi_n^2
 +\frac12(\boldsymbol\nabla\cdot\mathbf A_n)^2
 +\frac12m_n^2\mathbf A_n^2-\frac12m_n^2(A_n^0)^2
\\
&+\mathbf J_n\cdot\mathbf A_n-J_n^0A_n^0
 +A_n^0\boldsymbol\nabla\cdot\boldsymbol\Pi_n\biggr].
\end{align*}
This is still quadratic in $\boldsymbol\Pi_n$, with a field-independent (and
somewhat simpler) coefficient of the quadratic term, so the integral over the
$\boldsymbol\Pi_n$s can be done by just replacing $\boldsymbol\Pi_n$ with its
value at the stationary point of the functional
$\sum_n\int d^3x\,\boldsymbol\Pi_n\cdot\dot{\mathbf A}_n-H-\Delta H$:
\[
\boldsymbol\Pi_n=\dot{\mathbf A}_n+\boldsymbol\nabla A_n^0.
\]
With $\boldsymbol\Pi_n$ eliminated in this way,
$\sum_n\int d^3x\,\boldsymbol\Pi_n\cdot\dot{\mathbf A}_n-H-\Delta H$ is just
the Lorentz-invariant Lagrangian with which we started.

In order to take account of the possible need to introduce auxiliary fields
like $A_n^0$, from now on we shall write the path-integral formula after
elimination of the canonical conjugates in terms of fields $\psi_\ell$ that
include both canonical fields $q_n$ and auxiliary fields $c_r$:
\begin{align}
&\bra{\Omega;\mathrm{out}}T\left\{
 \mathcal O_A[\Psi_A(t_A)],\mathcal O_B[\Psi_B(t_B)],\ldots\right\}
 \ket{\Omega;\mathrm{in}}
\notag\\
&\quad\propto\int\prod_{\tau,\mathbf{x},n}d\psi_n(\mathbf{x},\tau)
 \mathcal O_A[\psi(t_A)]\mathcal O_B[\psi(t_B)]\cdots
\notag\\
&\qquad\cdot\exp\left[i\int_{-\infty}^{\infty}d\tau
 \left\{L[\psi(\tau),\dot\psi(\tau)]+i\epsilon\text{ terms}\right\}\right],
\tag{9.3.11}\label{eq:9.3.11}
\end{align}
it now being understood that $L$ includes any terms arising from a possible
field-dependent factor $(\operatorname{Det}\mathcal A)^{-1/2}$.
